\section{Conclusion}
To conclude, we presented the problem of building teams for funding that allows for collaboration opportunities. We then created and implemented  AI methods using string, taxonomy, and more advanced contextual boosted bandits in ULTRA, and demonstrated them to be quite useful both quantitatively (where informed methods increase recommendation quality while reducing their volume) and qualitatively in real human evaluations. 
We showed the generality of our approach in two different settings from US and India, and discussed our experience of deploying the system.


One area to extend our work is by considering larger data sizes for both researchers and RFPs, and from more diverse  sources. Furthermore, since our methods are dependent on open data about researchers as well as proposal calls, this dependency can be both a source of strength and weakness. Data has the potential to encourage  teaming without human bias but can also lead to inferior recommendation if the data is obsolete.  %if the data is obsolete or does not accurately reflect a researcher's  past or   future research directions, this can impact the features the teaming methods  extract, and therefore, the demanded skills may not match with researchers' backgrounds. 
Similarly, any feedback on the success of recommendation is only possible when a proposal has been won, and this data is usually not available or quite delayed (months or years after recommendation) to be useful. Considering a longer time frame for recommendation with proposal success data could lead to better results. Another area is to further enhance our goodness score with more metrics (e.g., considering the relevance of a researcher's previous projects to current RFPs and the number of ongoing projects a researcher is engaged in).

%ULTRA relies on publicly available profiles of researchers in order to make teaming suggestions. 
% However, this data may be obsolete or not accurately reflect their present and future research directions. This impacts the features the teaming methods have extracted, and therefore, the demanded skills may not match with researchers' backgrounds.

Our work inspires several interesting future extensions: considering a variety of domain knowledge including but not limited to fairness constraints, teaming constraints, domain constraints, etc. and developing a knowledge-driven learning system that can both exploit both the data and such knowledge remains an exciting future direction. A second direction would be to develop methods that would allow for interactive teaming where the system could not only present the recommendations but explain why and allow for human inputs to be used for refining the learned models. In addition, the scale of our survey could also be improved to ask about more aspects of the teams from participants: diversity and the connection strength between members, etc. A final and important direction is to scale this to different collaborative settings including but not limited to: healthcare, finance, law/legal, mental health, and educational support. 
